% ----------------------------------------
 \chapter{Conclusiones}
% ----------------------------------------
\label{C:conclusiones}

Se han estudiado las características correspondientes al microprocesador CDP1802 y entendido el principio de funcionamiento del mismo, conocimiento que fue aplicado para el desarrollo de diferentes circuitos de complejidad progresiva que se encuentran explicados en el presente documento. Se han reconocido los fundamentos y funcionalidades de los elementos correspondientes a una microcomputadora, como ser memorias, interfaz de entrada salida, conversor analógico digital, entre otras, de tal manera de poder proyectar subsistemas dedicados a funciones específicas que tengan interoperabilidad entre los mismos, seleccionando los diferentes dispositivos adecuadamente para su compatibilidad en las diferentes aplicaciones. Se analizaron circuitos y placas de desarrollo con el microprocesador utilizado, que sirvieron de base para los sistemas proyectados en este trabajo. 

Se han desarrollado competencias en el análisis de costos, pedidos de presupuesto y toda la logística involucrada en el proyecto e implementación de un circuito electrónico, comprendiendo las utilidades presentadas por un proceso de planificación y administración del tiempo, como el diagrama de Gantt, y también se ganó experiencia con herramientas de redacción de informes que permiten plasmar los resultados obtenidos en cualquier ámbito de manera profesional. 

Se han estudiado, analizado, comprendido y simulado diferentes arquitecturas que utilizan el microprocesador CDP1802, cuyos resultados brindaron entendimiento respecto a las topologías propuestas. Se utilizaron herramientas de simulación (Emma02, Proteus) que presentan herramientas de depuración, las cuales fueron sumamente importantes para detectar errores que presentan complejidad para detectarse en el circuito físico y corroborar el funcionamiento de las placas de desarrollo ya existentes con arquitecturas similares a las implementadas.

Se han proyectado, diseñado e implementado dos microcomputadoras (sistema mínimo y sistema medio) con arquitecturas diferentes y se ha logrado su correcto funcionamiento al ser probadas y testeadas. Simples programas se han implementado en el sistema mínimo y en el sistema medio, corroborando el correcto funcionamiento de los mismos. En el sistema medio se han ejecutado un intérprete de un lenguaje de programación y un sistema operativo sin disco, y se ha comprobado su correcto funcionamiento. Se ha proyectado y diseñado una microcomputadora (sistema final) que posea todas las aplicaciones planteadas desde un principio, como ser el conversor analógico digital y el uso de memorias seriales, por nombrar algunas. Se diseñaron el esquemático y el PCB del sistema final los cuales se encuentran presentes en el documento. Se logró integrar todas las funcionalidades pensadas en un principio y plasmarlas sobre el esquemático del sistema. 

Durante el desarrollo del trabajo se han utilizado y aplicado con éxito gran parte de los conocimientos obtenidos en la carrera de ingeniería electrónica y se adquirieron otros relacionados con llevar a la práctica los conceptos teóricos ya estudiados. Uno de ellos es el uso de software de diseño KiCad, en el cual se desarrollaron todos los circuitos planteados para el trabajo y la construcción de placas PCB de doble faz, entre otros. 


Como conclusión en vista de las facultades adquiridas con el desarrollo del presente trabajo, se destaca la interdependencia de software-hardware, sobre todo cuando el desarrollo de software es en bajo nivel como en este caso. Es de carácter fundamental para los diseñadores de este tipo de circuitos conocer el funcionamiento claro, uno del otro, para poder realizar un avance significativo en el desarrollo del proyecto. 

Las principales dificultades en el desarrollo del trabajo se relacionan con el entendimiento del microprocesador, el cual presenta una arquitectura particular. Se destaca principalmente el modo en el que opera con las instrucciones de entrada-salida, sin embargo, el uso de registros y su programación difieren mucho de las arquitecturas estudiadas en la carrera. 

Se han cumplimentado en gran medida los objetivos planteados para este proyecto en un principio, logrando integrar los conceptos estudiados a lo largo de la carrera correctamente y diseñar e implementar diferentes sistemas de microcómputo que sirven para sus determinadas aplicaciones. Los alumnos a cargo del proyecto han adquirido herramientas específicas para la integración de conceptos, proyecto y diseño de sistemas integrados, comprensión cabal del funcionamiento de diversos dispositivos, entre las más importantes. Estas herramientas, adquiridas por los alumnos para su próxima inserción laboral, cumplimentan así los objetivos de la materia Proyecto y Diseño Electrónico. 